\section{Short near-optimal scalar paths and residual spectral control}
\label{sec:critical-short-path-control}

Near-optimality controls the entropy spent at every prefix of an
increasing scalar path. This yields the path-length estimate used in
the upper bound, and controls the spectral loss when the path is lifted
to vector features. We use $V$, $V_*$ and the constant $c>0$ from
\eqref{eq:scalar-step-dissipation}.

\begin{lemma}[Dyadic bounds for a near-optimal increasing path]
\label{lem:critical-dyadic-short-path}
There is a universal $C$ such that, for every $0<p\le1/256$, an
increasing scalar path from $p$ with total cost at most $V(p)+p$ has
length $m$ and states $q_0=p,q_1,\ldots,q_m$ satisfying
\[
 m\le C/p,\qquad
 \sum_{j=0}^{m-1}q_j\le C\log(2/p).
\]
Such a path exists by the increasing-path reduction. The same sum bound
holds with the post-step states $q_1,\ldots,q_m$.
\end{lemma}

\begin{proof}
Write $D_j$ for the entropy cost of step $j$. Every suffix from $q_i$
including its readout costs at least $V(q_i)$, so
\begin{equation}\label{eq:critical-prefix-cost}
 \sum_{j=1}^{i}D_j
 \le V(p)+p-V(q_i)\le V_*-V(q_i)+p.
\end{equation}
Consider a dyadic band $[u,2u)$, where $u=2^kp<1/256$.
Because states increase, its transitions form one consecutive segment.
At most one transition exits the band. Every other transition has
post-state below $2u$, and the prefix ending at the last such transition
has cost, by~\eqref{eq:critical-prefix-cost}, at most
\[
 V_*-V(2u)+p\le34u+p\le35u.
\]
Every step whose pre-state lies in this band has cost at least $cu^2$,
by increasing states and strict dissipation. Thus at most
$35/(cu)+1$ steps start in the band, contributing at most $70/c+2u$
to the sum of their pre-states.

Choose the last band so that its next lower endpoint lies in
$[1/256,1/128)$. All remaining pre-states are at least $1/256$.
Since total cost is less than three, at most $3\cdot256^2/c$
steps remain. Summing the dyadic counts gives $m\le C/p$, and summing
the pre-state contributions gives $C\log(2/p)$.
Finally, the post-state sum differs from the pre-state sum by
$q_m-p\le1$.
\end{proof}

\begin{lemma}[A lower spectral bound for a lifted residual]
\label{lem:critical-residual-spectral}
Let $f\in\{-1,1\}^M$, let $R\succeq0$ have constant diagonal $r$,
and let $G\sim\mathcal N(0,R)$ be independent of a scalar random
variable $\xi$. Define vector features
\[
 h_x=\tanh(f_x\xi+G_x),\qquad
 \psi_r(z)=\mathbb E\tanh(z+\sqrt r Z),
\]
\[
 v'=\mathbb E\psi_r(\xi)^2,\qquad
 q'=\mathbb E\tanh^2(\xi+\sqrt r Z),\qquad
 R'=\mathbb Ehh^\top-v'ff^\top.
\]
Then
\begin{equation}\label{eq:critical-residual-spectral-bound}
 R'\succeq(1-q')^2R.
\end{equation}
This applies to arbitrary scalar control laws with finite second moment,
not only Gaussian controls.
\end{lemma}

\begin{proof}
Condition on $\xi$. The conditional mean of $h$ is $f\psi_r(\xi)$.
Put $a(\xi)=\mathbb E\operatorname{sech}^2(\xi+\sqrt rZ)$.
Gaussian integration by parts gives
\[
 \operatorname{Cov}(h,G\mid\xi)=a(\xi)R.
\]
The coefficient is the same at every input because the derivative is
even and the residual diagonal is constant. The covariance matrix of
$h-\mathbb E[h\mid\xi]-a(\xi)G$ is positive semidefinite, hence
\[
 \operatorname{Cov}(h\mid\xi)\succeq a(\xi)^2R.
\]
This argument also handles singular $R$ by Gaussian approximation, or
directly by a linear representation of $G$.
Average over $\xi$ and apply Jensen:
$R'=\mathbb E\operatorname{Cov}(h\mid\xi)
\succeq(\mathbb Ea(\xi))^2R=(1-q')^2R$.
\end{proof}

\begin{corollary}[Only polynomial spectral loss through a near-optimal lift]
\label{cor:critical-polynomial-residual-loss}
Lift a scalar path from Lemma~\ref{lem:critical-dyadic-short-path} by
the earlier scalar-lifting construction. Suppose the initial residual
has constant diagonal at most $\delta\le p$. Then for universal
$c_0>0,C_0<\infty$, its final residual satisfies
\[
 R_m\succeq c_0p^{C_0}R_0.
\]
\end{corollary}

\begin{proof}
At each step the scalar control is rescaled by the square root of the
ratio of the tracked component $v$ to the nominal state $q$.
The tracked component never exceeds the nominal state. To check this,
for $z\ge0$ and $r\ge0$, averaging the tanh addition formula at
opposite Gaussian shifts gives
\[
 0\le\psi_r(z)\le\tanh z.
\]
If $v\le q$, rescaling by $\sqrt{v/q}\le1$ can therefore only decrease
the absolute scalar feature. Induction from $v=q=p$ proves the claim.
The residual diagonal stays at most $\delta$ by conditional Gaussian
Poincare. Thus the actual total feature variance at step $j$ is at most
$q_j+\delta$.

Every nominal step has entropy at most three and reference variance
at most one. Gaussian entropy duality gives its preactivation second
moment at most $12+2\log2<14$. Rescaling does not increase this moment;
adding the independent residual raises it by at most one. Consequently
for each actual total post-step variance $\widehat q_j$,
\[
 1-\widehat q_j
 =\mathbb E\operatorname{sech}^2(\xi+\sqrt rZ)
 \ge\mathbb E e^{-2|\xi+\sqrt rZ|}
 \ge e^{-2\sqrt{15}}=:a_0>0.
\]
For $0\le u\le1-a_0$, one has $\log(1-u)\ge-u/a_0$.
Lemma~\ref{lem:critical-residual-spectral} therefore gives
\[
 R_m\succeq
 \exp\left[-\frac2{a_0}\sum_{j=1}^m\widehat q_j\right]R_0
 \succeq
 \exp\left[-\frac2{a_0}
       \left(\sum_{j=1}^mq_j+m\delta\right)\right]R_0.
\]
The dyadic lemma bounds the state sum by $C\log(2/p)$ and
$m\delta\le C$, proving the assertion.
\end{proof}

These are variational-path and population-covariance statements.
They do not alone give a finite-width lower probability or a sample
lower bound. Their role is to avoid an unnecessarily exponential
spectral loss when constructing such a finite-width comparison.
